%!TeX root=MemoriaTFG.tex

\chapter{Sparse Autoencoders y diccionario de conceptos clínicos}
\label{cap:anexof}

Este anexo documenta el módulo de interpretabilidad conceptual del
trabajo, basado en Sparse Autoencoders (SAE) sobre las
representaciones de PanDerm Large, su evaluación frente al
diccionario SkinCon~\cite{skincon} y la propuesta de extensión
con conceptos clínicos adicionales validados por la
Dra.~R.~Taberner. El módulo se incorpora al prototipo DermApIxel
como componente M3 (secci\'on~\ref{sec:dermapixel-m3}) y constituye una
contribución complementaria del trabajo
(secci\'on~\ref{sec:contrib-anexos}).

\section{Motivación}
\label{sec:anexof-motivacion}

La interpretabilidad de los modelos fundacionales constituye un
requisito creciente en escenarios clínicos donde la trazabilidad
del razonamiento es necesaria para la aceptación profesional. Los
encoders auto-supervisados como PanDerm producen representaciones
densas (vectores de $1\,024$ dimensiones para la variante
\emph{Large}) cuyos componentes individuales no admiten
interpretación directa. Los Sparse
Autoencoders~\cite{templeton2024} descomponen estas
representaciones en un diccionario disperso de \emph{features},
con la propiedad de que cada \emph{feature} puede asociarse, en
buena parte de los casos, a un concepto clínicamente
identificable. DermFM-Zero~\cite{dermfmzero2026} adopta este
mecanismo de interpretabilidad como contribución de su diseño,
pero sus pesos no son públicos (secci\'on~\ref{sec:lim-pesos}). El
presente trabajo replica el procedimiento sobre PanDerm Large con
pesos liberados y lo valida sobre SkinCon como ground-truth
conceptual.

\section{Arquitectura del SAE}
\label{sec:anexof-arquitectura}

La arquitectura del SAE empleado en el trabajo se compara con la
descrita por los autores de DermFM-Zero sobre PanDerm Base. La
tabla~\ref{tab:anexof-config} resume ambas configuraciones.

\begin{table}[H]
\centering
\footnotesize
\setlength{\tabcolsep}{4pt}
\renewcommand{\arraystretch}{0.95}
\caption{Configuración comparada del SAE entrenado en el trabajo
(``SAE Large'') y del SAE publicado por DermFM-Zero
(``SAE Base'').}
\label{tab:anexof-config}
\begin{tabular}{lll}
\hline
Parámetro                  & SAE Base (DermFM-Zero)             & SAE Large (este trabajo) \\
\hline
Encoder visual base        & PanDerm Base (ViT-B/16)            & PanDerm Large (ViT-L/16) \\
Dimensión de entrada       & $768$                              & $1\,024$ \\
\emph{Features} aprendidas & $6\,144$ ($8\times$)               & $16\,384$ ($16\times$) \\
Arquitectura               & TiedBias + UnitNormDecoder         & Pre/PostBias + Linear + UnitNormDecoder \\
Pérdida                    & $L_1$ ($3\times 10^{-5}$) + $L_2$  & $L_1$ ($1\times 10^{-3}$) + MSE \\
Datos de entrenamiento     & No especificado (pipeline DermFM-Zero) & $50\,454$ imágenes ($7$ datasets) \\
Épocas                     & $200$                              & $100$ (convergido) \\
Sparsity media observada   & ---                                & $16{,}5\%$ ($2\,753 \pm 112$ activas) \\
\hline
\end{tabular}
\end{table}

Los siete datasets de entrenamiento del SAE Large son HAM10000
($10\,015$ imágenes), BCN20000 ($12\,413$),
Dermnet ($19\,559$), PAD-UFES-20 ($2\,298$),
Fitzpatrick17k ($3\,887$), DDI ($647$) e HIBA
($1\,635$), con un total de $50\,454$ imágenes. La
composición exacta se documenta en el Anexo~\ref{cap:anexob}
(tabla~\ref{tab:anexob-sae-train}). La elección de
$L_1$ ($1\times 10^{-3}$) en lugar del valor publicado por
DermFM-Zero ($3\times 10^{-5}$) responde a la dimensión de entrada
mayor: el factor $L_1$ relativo a la magnitud típica del
\emph{embedding} se mantiene comparable.

\section{Comportamiento del SAE Large frente a SAE Base}
\label{sec:anexof-comparacion}

La tabla~\ref{tab:anexof-auroc} compara la AUROC de cuatro tareas
de discriminación binaria sobre los \emph{embeddings}
crudos (sin SAE) y sobre las activaciones del SAE para ambas
configuraciones. Las cifras corresponden a la media y desviación
estándar sobre veinte particiones aleatorias del conjunto de
evaluación.

\begin{table}[H]
\centering
\footnotesize
\setlength{\tabcolsep}{4pt}
\renewcommand{\arraystretch}{0.95}
\caption{Discriminación binaria sobre cuatro tareas con
\emph{embeddings} crudos y con activaciones del SAE. Media
$\pm$ desviación estándar sobre $20$ particiones aleatorias.}
\label{tab:anexof-auroc}
\begin{tabular}{lrrrr}
\hline
Tarea            & Raw $768$D (Base) & SAE Base ($6\,144$) & Raw $1\,024$D (Large) & SAE Large ($16\,384$) \\
\hline
Melanoma         & $0{,}906 \pm 0{,}020$ & $0{,}897 \pm 0{,}019$ & $0{,}897 \pm 0{,}020$ & $\mathbf{0{,}900 \pm 0{,}019}$ \\
Sesgo pelo       & $0{,}888 \pm 0{,}025$ & $0{,}879 \pm 0{,}023$ & $0{,}884 \pm 0{,}023$ & $\mathbf{0{,}889 \pm 0{,}024}$ \\
Sesgo tinta      & $\mathbf{0{,}848 \pm 0{,}024}$ & $0{,}836 \pm 0{,}025$ & $0{,}846 \pm 0{,}028$ & $0{,}851 \pm 0{,}028$ \\
Sesgo regla      & $0{,}888 \pm 0{,}024$ & $0{,}889 \pm 0{,}026$ & $0{,}894 \pm 0{,}022$ & $\mathbf{0{,}908 \pm 0{,}022}$ \\
\hline
\end{tabular}
\end{table}

El SAE Base degrada consistentemente la AUROC respecto a los
\emph{embeddings} crudos correspondientes (promedio
$-0{,}008$\,pp), mientras que el SAE Large produce mejoras de
$+0{,}007$\,pp de promedio sobre los \emph{embeddings} crudos
de $1\,024$ dimensiones. La interpretación operativa es que el
SAE Large preserva la información discriminativa relevante para
clasificación binaria pese a la proyección al espacio disperso de
$16\,384$ \emph{features}, comportamiento que justifica su
uso como módulo de interpretabilidad sin pérdida de utilidad
diagnóstica.

\begin{figure}[H]
\centering
\includegraphics[width=0.82\textwidth]{fig_sae_comparison.png}
\caption{Comparativa AUROC entre \emph{embeddings} crudos y
activaciones del SAE sobre cuatro tareas binarias. El SAE Large
(este trabajo, $1\,024 \to 16\,384$) preserva o mejora la
información discriminativa frente al SAE Base
($768 \to 6\,144$), que degrada consistentemente la AUROC.}
\label{fig:anexof-sae}
\end{figure}

\section{Concept Bottleneck Model sobre SkinCon}
\label{sec:anexof-cbm}

La validación clínica del SAE Large se realiza mediante un
\emph{Concept Bottleneck Model} (CBM) sobre el dataset
SkinCon~\cite{skincon}, que anota densamente $3\,230$
imágenes con $48$ conceptos clínicos por dermatólogos
certificados. La hipótesis a contrastar es que las
$16\,384$ \emph{features} del SAE Large contienen información
suficiente para predecir los conceptos SkinCon mediante un
clasificador lineal sobre cada concepto.

El procedimiento es el siguiente: cada imagen de SkinCon se codifica
mediante PanDerm Large; su \emph{embedding} de $1\,024$
dimensiones se proyecta al espacio disperso del SAE Large; sobre
las $16\,384$ activaciones se entrena un clasificador lineal
binario por concepto (regresión logística L-BFGS con
$C = 1{,}0$, \texttt{max\_iter}~$= 5\,000$,
\texttt{random\_state}~$=42$). La métrica reportada es la AUROC
del mejor \emph{feature} individual como predictor del concepto.

\subsection{Conceptos con mejor predictor}
\label{sec:anexof-cbm-top}

La tabla~\ref{tab:anexof-cbm} reporta los conceptos SkinCon con
mejor predictor entre las $16\,384$ \emph{features} del SAE
Large, ordenados por AUROC del mejor \emph{feature} descendente.

\begin{table}[H]
\centering
\footnotesize
\setlength{\tabcolsep}{4pt}
\renewcommand{\arraystretch}{0.95}
\caption{Conceptos SkinCon con mejor predictor entre las
\emph{features} del SAE Large. $n_+$: número de imágenes positivas
para el concepto en SkinCon. AUROC: del mejor \emph{feature}
individual del SAE como predictor binario del concepto.}
\label{tab:anexof-cbm}
\begin{tabular}{lrrr}
\hline
Concepto clínico             & $n_+$ & Mejor \emph{feature} & AUROC \\
\hline
Pedunculated                 & $26$  & \#1930 & $0{,}929$ \\
Friable                      & $153$ & \#5973 & $0{,}913$ \\
Exophytic/Fungating          & $42$  & \#10114 & $0{,}908$ \\
Comedo                       & $24$  & \#7191 & $0{,}903$ \\
Ulcer                        & $154$ & \#13908 & $0{,}897$ \\
Exudate                      & $144$ & \#5973 & $0{,}884$ \\
Wheal                        & $21$  & \#9552 & $0{,}884$ \\
Black                        & $90$  & \#3926 & $0{,}846$ \\
Telangiectasia               & $100$ & \#11172 & $0{,}828$ \\
Nodule                       & $189$ & \#1930 & $0{,}818$ \\
Umbilicated                  & $49$  & \#13783 & $0{,}810$ \\
\hline
\end{tabular}
\end{table}

Los conceptos con mejor predictor incluyen un núcleo de morfología
elevada (\emph{pedunculated}, \emph{exophytic/fungating},
\emph{nodule}: AUROC $0{,}818$--$0{,}929$), de
superficie alterada (\emph{friable}, \emph{ulcer}, \emph{exudate}:
AUROC $0{,}884$--$0{,}913$) y de estructura específica
(\emph{telangiectasia}: AUROC $0{,}828$). Cada \emph{feature}
del SAE puede actuar como predictor para más de un concepto
relacionado (e.g.~\#1930 cubre \emph{pedunculated} y
\emph{nodule}; \#5973 cubre \emph{friable} y \emph{exudate}), lo
que sugiere que el espacio disperso del SAE codifica dimensiones
clínicas de granularidad superior a la del concepto individual.

\subsection{Rendimiento agregado del CBM}
\label{sec:anexof-cbm-agregado}

Sobre los $22$ conceptos de SkinCon con cobertura
suficiente para el contraste, el protocolo \emph{CBM-Selected}
(regresión logística entrenada sobre $134$ \emph{features}
del SAE Large correlacionadas con cada concepto) alcanza AUROC
media de $0{,}880$. El protocolo de comparación
\emph{LP-Direct} (regresión logística sobre el \emph{embedding}
crudo de $1\,024$ dimensiones de PanDerm Large, sin pasar por
el SAE) alcanza AUROC media de $0{,}864$ sobre los mismos
conceptos. CBM-Selected supera a LP-Direct en $16$ de los
$22$ conceptos evaluados. El resultado indica que las
\emph{features} dispersas del SAE preservan ---y en la mayoría de
conceptos exceden--- el poder predictivo de la representación
densa subyacente, comportamiento que sostiene su uso como módulo
de interpretabilidad en el prototipo descrito en la
secci\'on~\ref{sec:dermapixel-m3} del Anexo~\ref{cap:anexoh}.

\begin{figure}[H]
\centering
\includegraphics[width=0.85\textwidth]{fig_skincon_cbm.png}
\caption{Concept Bottleneck Model sobre SkinCon: AUROC del mejor
\emph{feature} individual del SAE Large como predictor binario
por concepto. Once conceptos superan AUROC $0{,}80$, con
máximos en \emph{pedunculated} ($0{,}929$) y \emph{friable}
($0{,}913$).}
\label{fig:anexof-cbm}
\end{figure}

\section{Propuesta de extensión: quince conceptos clínicos adicionales}
\label{sec:anexof-propuesta}

El diccionario SkinCon, construido en el contexto del estudio de
Daneshjou et~al.~\cite{skincon}, prioriza conceptos visuales
estructurales de detección automatizable. La revisión clínica
realizada por la Dra.~R.~Taberner sobre el dataset
DermapixelAI~1.0 ha identificado $15$ conceptos clínicos
adicionales no contemplados por SkinCon que aparecen con
frecuencia relevante en la práctica dermatológica hispanohablante
y cuya inclusión enriquecería el diccionario operativo del
prototipo. La tabla~\ref{tab:anexof-propuesta} enumera los
conceptos propuestos.

\begin{table}[H]
\centering
\footnotesize
\setlength{\tabcolsep}{4pt}
\renewcommand{\arraystretch}{0.95}
\caption{Quince conceptos clínicos adicionales propuestos por la
Dra.~R.~Taberner sobre la base del diccionario SkinCon.}
\label{tab:anexof-propuesta}
\begin{tabular}{lll}
\hline
Categoría                       & Concepto propuesto                & Relevancia clínica \\
\hline
Pigmentación                    & Pigmentación reticular            & Diagnóstico diferencial nevus \\
Pigmentación                    & Pseudored                         & Lentigo solar \emph{vs.}\ lentigo maligno \\
Vasculatura                     & Lagunas vasculares                & Hemangioma, angioqueratoma \\
Vasculatura                     & Vasos en horquilla                 & Carcinoma basocelular \\
Estructura epidérmica           & Hiperqueratosis                    & Queratosis seborreica \\
Estructura epidérmica           & Comedo invertido                   & Queratosis seborreica adelgazada \\
Patrón inflamatorio             & Eritema heliotropo                 & Dermatomiositis \\
Patrón inflamatorio             & Patrón en \emph{papel de seda}     & Liquen escleroso \\
Distribución                    & Distribución dermatomérica         & Herpes zóster \\
Distribución                    & Distribución \emph{en alas}        & Lupus subagudo \\
Superficie                      & Cráteres centrales                 & Molluscum contagioso \\
Superficie                      & Anillos concéntricos               & Eritema multiforme \\
Color                           & Coloración \emph{salmón}           & Psoriasis guttata \\
Color                           & Halo blanco perilesional           & Nevus halo \\
Otros                           & Bordes geográficos                 & Vitíligo segmentario \\
\hline
\end{tabular}
\end{table}

La incorporación de estos quince conceptos al diccionario operativo
del prototipo DermApIxel constituye una línea de continuación
inmediata, condicionada al muestreo de aproximadamente
$200$ imágenes anotadas por concepto para entrenar el
clasificador lineal correspondiente sobre las activaciones del SAE
Large. El protocolo de anotación previsto utiliza la interfaz
autocontenida descrita en la secci\'on~\ref{sec:lineas-abiertas} del
capítulo~\ref{cap:conclusiones}.

\section{Limitaciones del módulo SAE}
\label{sec:anexof-limitaciones}

El módulo SAE presenta tres limitaciones declaradas. Primero, la
asociación entre \emph{feature} y concepto se establece sobre la
base del mejor predictor individual, sin garantía de que la
\emph{feature} sea \emph{monosemántica} en el sentido estricto del
término (la misma \emph{feature} puede correlacionarse con varios
conceptos relacionados, como se observa en
secci\'on~\ref{sec:anexof-cbm-top}). Segundo, el rendimiento del SAE Large
sobre tareas de discriminación binaria
(tabla~\ref{tab:anexof-auroc}) es ligeramente superior al de los
\emph{embeddings} crudos correspondientes, pero esta mejora es del
orden de la desviación estándar inter-partición, lo que limita la
afirmación de \emph{superioridad estadísticamente significativa}
sin las pruebas pendientes (secci\'on~\ref{sec:lim-estadistica}).
Tercero, la validación sobre conceptos clínicos se realiza sobre
SkinCon, dataset con sesgo de selección hacia patología
dermatoscópica documentada, lo que condiciona la transferencia a
otros contextos clínicos.
